\section*{Chapitre \ref{ch:m1:matching-algorithms-and-implied-spreads} --- Algorithmes d'appariement et spreads implicites}

\begin{solution}{exo:m1:matching-algorithms-and-implied-spreads:1}
A 10, B 200, C 20 sur ses 40, D rien.
\end{solution}

\begin{solution}{exo:m1:matching-algorithms-and-implied-spreads:2}
Actions $60 \times (10, 200, 40, 250)/500 = 1.2$, 24, 4.8, 30. Parties entières 1, 24, 4,
30~: 59 lots. Le lot restant revient à l'ordre le plus ancien, A. Exécutions~: A 2, B
24, C 4, D 30.
\end{solution}

\begin{solution}{exo:m1:matching-algorithms-and-implied-spreads:3}
Cours acheteur $10\,050 - 10\,023 = 27$ pour $\min(30, 10) = 10$. Cours vendeur $10\,052 - 10\,020 =
32$ pour $\min(25, 40) = 25$.
\end{solution}

\begin{solution}{exo:m1:matching-algorithms-and-implied-spreads:4}
Réduire la taille conserve la place. L'augmenter, changer le prix, et
annuler puis ressaisir la font perdre (le dernier cas par construction). Pour ajouter
de la taille sans perdre la priorité, il ne faut jamais modifier à la hausse~: on laisse l'ancien ordre à sa place
et saisir un \emph{second} ordre pour la quantité supplémentaire, qui
se place à la fin de la file. Le gestionnaire d'ordres doit donc gérer plusieurs
ordres actifs par niveau de prix et annuler en premier le plus récent lorsqu'il réduit.
\end{solution}

\begin{solution}{exo:m1:matching-algorithms-and-implied-spreads:5}
$8 \times 1\,000/92 = 86.96$, soit 87~; $\lfloor 100 \times 87/1\,087\rfloor = 8$,
$\lfloor 100 \times 86/1\,086\rfloor = 7$, $\lfloor 1\,000 \times 87/1\,087\rfloor
= 80$. Si chacun multiplie par dix, les parts ne changent pas~: chacun reçoit
ce qu'il recevait auparavant, la profondeur affichée est dix fois plus grande, et chacun est
exposé à dix fois plus de volume lors d'un balayage. La profondeur affichée mesure alors la
tolérance des participants à être surexécutés, et non leur volonté d'échanger, ce qui
explique que la profondeur des produits à \omterm{def:m1:matching-algorithms-and-implied-spreads:prorata}{allocation au prorata} s'évapore dans les marchés rapides.
\end{solution}

\begin{solution}{exo:m1:matching-algorithms-and-implied-spreads:6}
Un ordre de 15 lots a été annulé (une transaction aurait produit un
message de transaction). S'il était devant l'ordre de la stratégie, 15 lots de moins le précèdent~; s'il était derrière,
rien n'a changé. Le flux par prix ne dit pas lequel. Une estimation courante
suppose que les annulations se répartissent proportionnellement à la quantité
devant et derrière, ou, plus prudemment, qu'elles viennent toutes de derrière. Le flux par
ordre donne l'identifiant de l'ordre annulé et donc sa place~: la
position est connue exactement.
\end{solution}

\begin{solution}{exo:m1:matching-algorithms-and-implied-spreads:7}
En tête d'une file FIFO~: 9.56 des 10 lots. Au prorata~: 2.34 quelle que soit
la position. Rapport 4.1. L'\omterm{def:m1:matching-algorithms-and-implied-spreads:prorata}{allocation au prorata} est meilleure pour un ordre ayant 170 lots ou plus
devant lui dans un niveau de 500 lots.
\end{solution}

\begin{solution}{exo:m1:matching-algorithms-and-implied-spreads:8}
La part exacte d'un agresseur de 100 lots est $100 \times 5/2\,005 = 0.249$ lot~: elle
s'arrondit à zéro, et avec un minimum de deux lots, aucun ordre de moins de 41
lots affiché ($\lfloor 100 \times 41/2\,041\rfloor = 2$) ne reçoit quoi que ce soit
au prorata. Les restes vont aux ordres les plus anciens. Le trader ne récolte rien
du flux ordinaire et n'est exécuté que lorsqu'un ordre de 800 lots ou plus
balaie le niveau~: un portefeuille composé exclusivement de mauvaises exécutions. Dans un produit à \omterm{def:m1:matching-algorithms-and-implied-spreads:prorata}{allocation au prorata}
il existe une taille minimale viable, et elle est grande.
\end{solution}

\begin{solution}{pb:m1:matching-algorithms-and-implied-spreads:1}
\textbf{1.} $10\,050 - 10\,023 = 27$ pour 10.
\textbf{2.} $10\,052 - 10\,020 = 32$ pour 25.
\textbf{3.} Cours acheteur 28 pour 5 (direct)~; cours vendeur 32 pour 32 (7 directs et 25 implicites).
\textbf{4.} Cela signifierait que le meilleur cours acheteur de l'échéance proche dépasse le
meilleur cours vendeur de l'échéance éloignée augmenté d'un prix auquel quelqu'un accepte de
vendre le spread~: le moteur apparierait les trois ordres d'un coup. Les implicites croisés
avec des ordres directs s'exécutent immédiatement~; ils ne sont jamais affichés.
\textbf{5.} Un cours acheteur implicite et un cours vendeur implicite qui se croisent seraient
composés uniquement d'ordres fermes qui ne se croisent pas dans leurs propres carnets~;
les apparier exigerait des chaînes de jambes sur d'autres instruments, sans limite.
Les bourses coupent la récursion~: chaque appariement contient au moins un ordre
direct.
\textbf{6.} Le cours acheteur direct, 5 à 28~; puis le cours acheteur implicite, 7 à 27.
\textbf{7.} $(5 \times 28 + 7 \times 27)/12 = 27.42$.
\textbf{8.} Pour la partie implicite~: 7 Z achetés par l'acheteur de l'échéance proche à
10\,050 et 7 H vendus par le vendeur de l'échéance éloignée à 10\,023. La partie directe
n'apparaît que dans le spread (les jambes reçoivent des prix attribués selon une
convention de la bourse et ne touchent pas aux carnets fermes).
\textbf{9.} Échéance proche~: cours acheteur 10\,050 $\times$ 23, cours vendeur inchangé. Échéance éloignée~: cours acheteur
inchangé, cours vendeur 10\,023 $\times$ 3. Spread~: cours acheteur 27 $\times$ 3 (implicite), cours vendeur 32
$\times$ 32.
\textbf{10.} Vendre 12 Z à 10\,050~; acheter 10 H à 10\,023 et 2 au
cours vendeur suivant, à 10\,024 au mieux~: un spread d'au plus 26.83, contre 27.42, avec le
risque que le marché bouge entre les deux jambes.
\textbf{11.} Cours acheteur $28 + 10\,020 = 10\,048$ pour 5~; cours vendeur $32 + 10\,023 = 10\,055$ pour
7.
\textbf{12.} Non~: l'échéance proche affiche 10\,050 à 10\,052, meilleur des deux
côtés.
\textbf{13.} $10\,050 - 10\,021 = 29$ pour $\min(30, 20) = 20$.
\textbf{14.} Non. Dans son propre carnet, le meilleur cours acheteur est à 10\,020, et le
cours acheteur \omterm{def:m1:matching-algorithms-and-implied-spreads:implied}{implied-out} de l'échéance éloignée, $b_F - a_S = 10\,050 - 32 = 10\,018$, est encore plus bas. Il s'exécuterait
si quelqu'un vendait le spread à 29 ou mieux, ou achetait à 10\,021 sur l'échéance éloignée.
\textbf{15.} Un tick derrière~: le meilleur cours acheteur affiché est désormais 29 pour 20,
implicite. Pour repasser en tête, le \omterm{def:m1:what-a-trading-firm-does:market-maker}{teneur de marché} doit acheter à 29, et la
position de cet ordre par rapport aux 20 lots implicites dépend de la règle du moteur.
Un ordre sur l'échéance éloignée vient de refixer le spread sans qu'aucun ordre de spread
n'ait été saisi~: les cotations d'un carnet doivent être pilotées par les trois.
\textbf{16.} Une telle approche manque les exécutions issues des vendeurs de spread par
des ordres \omterm{def:m1:matching-algorithms-and-implied-spreads:implied}{implied-out}, et elle se trompe sur la position dans la file~: la quantité \omterm{def:m1:matching-algorithms-and-implied-spreads:implied}{implied-out}
peut se placer devant ou derrière l'ordre de la stratégie au même prix. Elle
ignore aussi que la profondeur visible de l'échéance proche comprend une quantité qui disparaît
lorsque l'échéance \emph{éloignée} bouge.
\textbf{17.} Si les implicites se classent derrière tous les ordres directs, la position
d'un ordre direct dans la file est protégée~; s'ils se classent selon l'heure de création, la quantité implicite
peut passer devant chaque fois qu'un autre carnet bouge. L'exécution espérée, et
sa \omterm{def:m1:what-a-trading-firm-does:adverse-selection}{sélection adverse}, diffèrent entre les deux cas.
\textbf{18.} Calculer et diffuser la profondeur implicite à chaque niveau de
chaque combinaison est coûteux, et elle change à chaque mise à jour de n'importe quelle jambe~;
le meilleur niveau est ce qui compte pour l'appariement.
\textbf{19.} \textbf{Meilleur cours acheteur implicite 27 pour 10~; meilleur cours acheteur affiché 28 pour 5.}
\textbf{20.} Il permet à un roll saisi comme un seul ordre de puiser dans la liquidité
des deux échéances à la fois, sans risque d'être exécuté sur une seule jambe seulement.
\end{solution}

\begin{solution}{iq:m1:matching-algorithms-and-implied-spreads:1}
Le FIFO exécute les ordres au repos au meilleur prix, du plus ancien au plus récent~; l'\omterm{def:m1:matching-algorithms-and-implied-spreads:prorata}{allocation au prorata} répartit
l'ordre entrant proportionnellement à sa taille. Les produits qui évoluent activement, avec
de petits ticks par rapport à leur volatilité (indices actions, pétrole brut, note futures) sont
FIFO~: la priorité s'acquiert en améliorant le prix ou en arrivant en premier.
Les produits qui bougent à peine par rapport à leur tick, en particulier les contrats à terme sur taux
à court terme, recourent à l'\omterm{def:m1:matching-algorithms-and-implied-spreads:prorata}{allocation au prorata} ou à des hybrides combinant \omterm{def:m1:matching-algorithms-and-implied-spreads:top}{allocation au premier ordre} et
allocations aux \omterm{def:m1:what-a-trading-firm-does:market-maker}{teneurs de marché}, de sorte que de nombreuses firmes continuent de coter à un prix qui
reste stable pendant des heures.
\iqlookfor{le lien entre la taille du tick par rapport à la volatilité et la règle.}
\end{solution}

\begin{solution}{iq:m1:matching-algorithms-and-implied-spreads:2}
Un ordre que le moteur d'appariement dérive d'ordres de carnets liés~: deux
ordres directs sur des échéances différentes impliquent un ordre de
\omterm{def:m1:matching-algorithms-and-implied-spreads:calendar}{calendar spread} (\omterm{def:m1:matching-algorithms-and-implied-spreads:implied}{implied-in})~; un ordre de spread et un ordre direct impliquent un ordre sur l'autre jambe
(\omterm{def:m1:matching-algorithms-and-implied-spreads:implied}{implied-out}). Il est négociable, toutes les jambes s'exécutent atomiquement, et un implicite
ne s'apparie jamais avec un implicite.
\iqlookfor{les deux sens, exécution atomique, aucun implicite contre implicite.}
\end{solution}

\begin{solution}{iq:m1:matching-algorithms-and-implied-spreads:3}
FIFO~: investir dans la latence pour atteindre en premier les nouveaux niveaux de prix, maintenir les ordres
en vie pour préserver sa position dans la file, évaluer chaque position dans la file, afficher la taille réelle.
\omterm{def:m1:matching-algorithms-and-implied-spreads:prorata}{Allocation au prorata}~: la vitesse compte moins~; afficher plus que prévu pour obtenir une part,
gérer le risque d'être sur-exécuté lors des balayages, annuler rapidement lorsque de gros
ordres apparaissent, et tenir compte des seuils minimaux d'allocation ainsi que de tout
privilège de l'\omterm{def:m1:matching-algorithms-and-implied-spreads:top}{allocation au premier ordre}, qui ramène la course à la vitesse pour le premier ordre
à un nouveau prix.
\iqlookfor{le surdimensionnement et son risque~; l'étape de l'\omterm{def:m1:matching-algorithms-and-implied-spreads:top}{allocation au premier ordre} qui réintroduit la vitesse.}
\end{solution}

\begin{solution}{iq:m1:matching-algorithms-and-implied-spreads:4}
Calculer $\lfloor a \cdot s_i / S\rfloor$ en nombres entiers pour chaque ordre, sur les quantités ouvertes~;
écarter les allocations inférieures au minimum~; répartir le reliquat selon
la règle secondaire. Cas limites~: agresseur plus grand que le niveau~; total des parties entières inférieur à l'agresseur~;
toutes les allocations inférieures au minimum (tout passe alors par la règle secondaire)~;
un ordre déjà partiellement exécuté lors d'une
étape précédente~; dépassement de $a \cdot s_i$~; ne jamais dépasser la taille d'un ordre~;
déterminisme du départage. Tester la conservation sur des carnets aléatoires.
\iqlookfor{arithmétique entière, conservation, et le cas où toutes les allocations sont inférieures au minimum.}
\end{solution}

\begin{solution}{iq:m1:matching-algorithms-and-implied-spreads:5}
Noter la quantité du niveau au moment où l'ordre est confirmé~: c'est la
quantité devant soi. Les transactions au niveau la réduisent unité pour unité. Les baisses sans
transaction sont des annulations~: les répartir devant et derrière selon une hypothèse
(proportionnelle, ou tout derrière pour une estimation prudente), bornée de sorte que
la quantité devant ne dépasse jamais la quantité du niveau moins son propre ordre.
Calibrer l'hypothèse sur ses propres exécutions~: les délais d'exécution réalisés
révèlent la position réelle a posteriori.
\iqlookfor{la borne issue du total du niveau, et la calibration sur ses propres exécutions.}
\end{solution}

\begin{solution}{iq:m1:matching-algorithms-and-implied-spreads:6}
Un moteur unique couvre toutes les jambes, puisqu'une mise à jour dans n'importe quel carnet peut créer ou
retirer de la quantité négociable dans les autres~; exécution atomique multi-jambes avec des prix
par jambe attribués selon la convention de la bourse~; la règle classant les implicites
face aux ordres directs à prix égal~; aucun implicite contre implicite~; des ticks différents
pour les spreads et les outrights~; l'algorithme d'allocation de chaque carnet~; et
des données de marché qui publient la quantité implicite de la même manière que le flux réel,
sinon les stratégies entraînées sur le simulateur verront un carnet qui
n'existe pas.
\iqlookfor{propagation des événements entre carnets et règle de classement.}
\end{solution}
